\chapter{Conclusiones y trabajo futuro}\label{cap:conclusiones}

Este capítulo cierra la memoria. Valora el grado de cumplimiento de los objetivos, resume
las aportaciones del trabajo, declara sus limitaciones sin suavizarlas y esboza las líneas de
trabajo futuro que el sistema habilita, con la comparación de recomendadores como principal
de ellas.

\section{Grado de cumplimiento de objetivos}\label{sec:conc-objetivos}

De los siete objetivos específicos fijados en la Sección~\ref{sec:intro-objetivos-especificos},
los cuatro de producto e infraestructura de datos están cumplidos y verificados. El catálogo
tiene identidad canónica propia, procedencia por obra y licencia documentada, con más de
trescientas mil obras servidas (OE-1). Una persona usuaria puede mantener su lista de
pendientes, valorar juegos e inventariar copias físicas y digitales con su formato y su
edición (OE-2). La búsqueda es tolerante. El catálogo se puede ordenar y filtrar por
plataforma, género, año y nota (OE-3). Existe una página de recomendaciones por afinidad de
géneros basada en la actividad propia de cada persona, visiblemente distinta del baseline de
popularidad (OE-4). El contrato de evaluación ya está congelado, y también se han
implementado y probado sus métricas, particiones y primeras primitivas de contenido (OE-5,
parcialmente). La comparación completa de recomendadores y el panel de investigación siguen
como trabajo futuro (OE-6 y OE-7).

Cada incremento del sistema se cerró con sus criterios de éxito verificados de forma
independiente y con la firma de aceptación del autor. El objetivo general se cumple en su
primera mitad y deja avanzada la segunda: la plataforma de catalogación existe, funciona y
sirve de base al banco de pruebas; el corpus gobernado, los snapshots de ratings y el
contrato del laboratorio ya están implementados. Falta ejecutar la población sintética
completa, las variantes finales y la comparación que los explota.

\section{Aportaciones}\label{sec:conc-aportaciones}

El trabajo deja tres aportaciones. La primera es el producto: una plataforma de
catalogación de videojuegos que combina inventario de propiedad física y digital, catálogo a
escala real y procedencia de datos documentada, un conjunto que los productos de referencia
no reúnen. La segunda es un banco de pruebas reproducible: un corpus citable, un modelo de
dominio pensado para la evaluación, usuarios simulados y una disciplina de procedencia sobre
la que se pueden montar el contrato de evaluación y los experimentos. La tercera es la
metodología de ingeniería asistida por agentes, aplicada y evaluada a lo largo de todo el
proyecto y versionada en el repositorio, con un árbol de planificación vivo, registros de
decisión, un registro de agentes de solo anexado y gates de evidencia, de modo que el proceso
que produjo el sistema es tan inspeccionable y reproducible por un tercero como el sistema
mismo.

\section{Limitaciones}\label{sec:conc-limitaciones}

Las limitaciones conocidas se declaran sin ocultarlas, como recogen los ítems diferidos del
repositorio. La Tabla~\ref{tab:limitaciones} las resume.

\begin{table}[htp]
\centering
\begin{tabularx}{\textwidth}{ | Y | Y | }
    \hline
    \textbf{Limitación} & \textbf{Situación} \\
    \hline
    Cobertura parcial de ratings & El rating de usuarios de IGDB está disponible en 27.014 obras gobernadas (13,9330\%); RAWG no añadió cobertura nueva en el corte de 10.000 obras. \\
    \hline
    Fecha de lanzamiento parcial & Poblada para en torno al 75\% de las filas. \\
    \hline
    Comparación de laboratorio no ejecutada & El protocolo, las métricas y las primeras primitivas están probados, pero todavía no existe el artefacto final con resultados de Precision, Recall, nDCG y MAP. \\
    \hline
    Pasada de refinamiento de diseño pendiente & Cuatro elementos cosméticos aceptados por el autor: mezcla de idiomas en la explicación de recomendaciones, navegación superior autenticada estrecha por debajo de 430 píxeles, detalle de juego algo escaso y menú de cuenta que no identifica la cuenta activa. \\
    \hline
    Accesibilidad manual del recorrido inicial & El reflujo al 400\% de zoom y una pasada con lector de pantalla siguen abiertos. \\
    \hline
    Despliegue no conectado a datos reales & Pospuesto al final del proyecto por decisión del autor; la demostración se ejecuta en local. \\
    \hline
\end{tabularx}
\caption{Limitaciones conocidas del sistema.}
\label{tab:limitaciones}
\end{table}

\section{Trabajo futuro}\label{sec:conc-trabajo-futuro}

El banco de pruebas reproducible que esta memoria deja construido habilita de forma natural
la comparación sistemática de recomendadores. La Fase 2 ya ha congelado un protocolo
consciente de la simulación: deja una interacción fuera por persona usuaria, fija una rejilla
de 18 configuraciones sobre validación, reserva el test para una única ejecución y establece
P@K, R@K, nDCG@K y MAP@K en K de 5, 10 y 20. También deja especificados el baseline
aleatorio, el de popularidad y las tres variantes de contenido, con explicaciones
deterministas y una vía de arranque en frío.

El siguiente trabajo no consiste en cambiar esas reglas al ver los resultados. Los
arquetipos sintéticos ya se han generado y validado con 200 perfiles reproducibles y una
cohorte cold-start separada. Queda cerrar las variantes y ejecutar el runner sobre el mismo
manifiesto de candidatos. Después se podrán añadir el filtrado colaborativo y los
híbridos, siempre como extensiones comparables y con nuevas versiones del protocolo si
introducen otras señales. El panel accesible deberá leer los artefactos producidos, exportar
sus tablas y declarar que las cifras son evidencia de simulación, no de usuarios reales.
Más allá de la recomendación, quedan los flujos completos de colección, la portabilidad de
datos y el descubrimiento público.

\section{Cierre}\label{sec:conc-cierre}

SavePoint parte de una idea doble: un producto de catalogación que a la vez sea un banco de
pruebas de recomendadores. El trabajo deja el producto construido, verificado y documentado,
con el banco de pruebas cimentado sobre datos citables y un modelo pensado para la
evaluación. Deja además una forma de trabajar, asistida por agentes y gobernada por
especificación, cuyo rastro completo de decisiones queda en el repositorio. La comparación de
recomendadores, que es el siguiente paso, tiene ahora una base sobre la que apoyarse.
